% Part: sets-functions-relations
% Chapter: sets
% Section: basics

\documentclass[../../../include/open-logic-section]{subfiles}

\begin{document}

\olfileid[id]{sfr}{set}{bas}
\olsection{Ekstensionalitas}

Sebuah \emph{himpunan} adalah kumpulan objek yang dipandang sebagai
satu objek. Objek-objek yang membentuk himpunan disebut \emph{elemen}
atau \emph{anggota} himpunan tersebut. Jika $x$ merupakan !!a{element}
himpunan~$A$, kita tulis $x \in A$; jika bukan, kita tulis $x \notin
A$. Himpunan yang tidak mempunyai !!{element}s disebut \emph{himpunan
kosong} dan dilambangkan dengan~``$\emptyset$''.

\begin{explain}
Tidak menjadi soal bagaimana kita \emph{menentukan} himpunan itu,
bagaimana kita \emph{mengurutkan} !!{element}s, atau bahkan berapa
\emph{kali} kita menghitung !!{element}s itu. Yang penting hanyalah
objek mana yang menjadi !!{element}s himpunan itu. Hal ini kita rumuskan
dalam prinsip berikut.
\end{explain}

\begin{defn}[Ekstensionalitas]
  Jika $A$ dan $B$ adalah himpunan, maka $A = B$ jika dan hanya jika
  setiap !!{element} dari~$A$ juga merupakan !!a{element} dari~$B$,
  dan sebaliknya.
\end{defn}

Ekstensionalitas membenarkan penggunaan suatu notasi. Secara umum,
jika kita mempunyai objek $a_{1}$, \dots, $a_{n}$, maka
$\{a_{1}, \dots, a_{n}\}$ adalah \emph{satu-satunya} himpunan dengan
$a_1, \ldots, a_n$ sebagai !!{element}s. Kita menekankan kata
``\emph{satu-satunya}'', sebab ekstensionalitas menyatakan bahwa hanya
ada \emph{satu} himpunan semacam itu. Ekstensionalitas juga membenarkan
kesamaan berikut:
  \[
    \{a, a, b\} = \{a, b\} = \{b,a\}.
  \]
Jadi, ketika membahas himpunan, kita tidak memedulikan urutan
!!{element}s ataupun berapa kali anggota itu dicantumkan.

\begin{tagblock}{novice}
\begin{ex}
Setiap kali kita mempunyai sekumpulan objek, kita dapat menghimpunnya
menjadi sebuah himpunan. Sebagai contoh, himpunan saudara kandung
Richard memuat satu orang dan dapat kita tulis sebagai
$S=\{\textrm{Ruth}\}$. Himpunan bilangan bulat positif yang kurang dari
$4$ adalah $\{1, 2, 3\}$, tetapi dapat pula ditulis $\{3, 2, 1\}$ atau
bahkan $\{1, 2, 1, 2, 3\}$. Berdasarkan ekstensionalitas, semua ini
adalah himpunan yang sama. Setiap !!{element} dari $\{1, 2, 3\}$ juga
merupakan !!a{element} dari $\{3, 2, 1\}$ (dan dari $\{1, 2, 1, 2,
3\}$), dan sebaliknya.
\end{ex}
\end{tagblock}

Sering kali kita menentukan sebuah himpunan melalui suatu sifat yang
dimiliki bersama oleh !!{element}s. Kita akan menggunakan
notasi ringkas $\Setabs{x}{\phi(x)}$, dengan $\phi(x)$ menyatakan sifat
yang harus dimiliki~$x$ agar dihitung sebagai salah satu !!{element}s
himpunan itu.

\begin{tagblock}{novice}
\begin{ex}
Dalam contoh kita, $S$ dapat pula ditentukan sebagai
\[
S = \Setabs{x}{x \text{ adalah saudara kandung Richard}}.
\]
\end{ex}
\end{tagblock}

\begin{tagblock}{math}
\begin{ex}
Suatu bilangan disebut \emph{sempurna} jika dan hanya jika bilangan itu
sama dengan jumlah pembagi sejatinya (yaitu bilangan-bilangan yang
membaginya sampai habis, tetapi tidak sama dengan bilangan tersebut).
Sebagai contoh, $6$ sempurna karena pembagi sejatinya adalah $1$, $2$,
dan~$3$, serta $6 = 1 + 2 + 3$. Bahkan, $6$ adalah satu-satunya
bilangan bulat positif kurang dari $10$ yang sempurna. Karena itu,
dengan menggunakan ekstensionalitas, kita dapat mengatakan:
\[
	\{6\} = \Setabs{x}{x\text{ sempurna dan }0 \leq x \leq 10}
\]
Notasi di sebelah kanan dibaca ``himpunan semua $x$ sedemikian sehingga
$x$ sempurna dan $0 \leq x \leq 10$''. Identitas ini menegaskan bahwa,
ketika membahas himpunan, kita tidak memedulikan cara himpunan tersebut
ditentukan. Secara lebih umum, ekstensionalitas menjamin bahwa selalu
hanya ada satu himpunan semua $x$ sedemikian sehingga $\phi(x)$.
Jadi, ekstensionalitas membenarkan penyebutan
$\Setabs{x}{\phi(x)}$ sebagai \emph{himpunan} semua $x$ sedemikian
sehingga~$\phi(x)$.
\end{ex}
\end{tagblock}

Ekstensionalitas memberi kita cara untuk menunjukkan bahwa dua
himpunan identik: untuk menunjukkan $A = B$, tunjukkan bahwa jika
$x \in A$, maka $x \in B$, dan jika $y \in B$, maka $y \in A$.

\begin{prob}
Buktikan bahwa terdapat paling banyak satu himpunan kosong, yakni
tunjukkan bahwa jika $A$ dan $B$ adalah himpunan tanpa !!{element}s,
maka $A = B$.
\end{prob}

\end{document}
